% Edition 21, section 13 initial segment. French physical pages 74--77.
% English physical pages 49--51, stopping before the proof of Lemma 13.3.
% Parent integrates this file; no inherited source is overwritten.
\egasection{§13\quad 初等态射与 M. Artin 定理}
\source{法文手稿\texttt{EGA5\_francais.pdf}物理第74--77页；英文编辑稿\texttt{EGAVaiello.pdf}物理及印刷第49--51页，第51页首段止}

% V21-INI-P01: French 74, definition 13.1.
\egas{13.1}{定义}预概形的态射$f\colon X\to Y$称为一个“初等态射”，如果$X$在$Y$上同构于形如$X'\setminus Z$的预概形，其中$X'\to Y$是光滑射影态射，其几何纤维连通且维数为1，而$Z$是$X'$的闭子预概形，使得态射$Z\to Y$平展、满射且次数恒定。\cmt{V21-INI001}{法稿第74页在定义中确实写成$X'\to X$，英文第49页已经改为$X'\to Y$。前面的“$Y$-同构”、边界态射$Z\to Y$以及整个定义均确定此处基底为$Y$，故正文保留正确的$X'\to Y$。这个错误属于法稿，不能反向归于英文。}

一个态射称为多重初等的，如果它是初等态射的复合。域$k$上的预概形$X$称为（在$k$上）多重初等的，如果其结构态射$X\to\Spec(k)$是多重初等的。\footnote{\textbf{英译编者注〔英文脚注63〕：}提醒Blass：向Artin询问术语。参见Milne第117页及SGA IV〔英译者注〕。}\cmt{V21-INI002}{“初等”“多重初等”沿用已有译稿。法稿定义中的旧术语可见删划，删去的字没有重新作为有效定义；不能可靠辨认的覆打字不予恢复。英文脚注63没有给出Milne的书名及SGA IV的具体位置，故不补造书目。关于允许空复合的说明实际出现在下面13.2的证明中，仍在那里保留作者的编写提醒。}

% V21-INI-P02: French 74, theorem 13.2 and first reduction.
\egas{13.2}{定理（M. Artin）}设$k$是无限完美域，$X$是$k$上的几何不可约预概形，$x$是$X$的一个光滑点。则$x$具有一个由多重初等开邻域组成的基本邻域系。\cmt{V21-INI003}{法稿第74页“几何不可约”“完美且无限”是明确的手添条件，均已保留；原陈述没有把$x$限为$k$-有理点。所需有限型归约可以从$x$处光滑性得到：在任意给定邻域内先取一个包含$x$的光滑仿射邻域，它在$k$上有限表示。详见本版核查报告“初段补证一”。原后续证明先使用有理点处的线性截面；一般闭点及任意光滑点的衔接，见报告对13.4(c)有限点构造的译者补证，不能把原稿的略写说成已逐项给出了该论证。}

把$X$替换为$x$的一个给定邻域后，只须证明$x$在$X$中有一个多重初等开邻域。\cmt{V21-INI004}{法、英两稿在这一句都写“初等开邻域”。所要证明的是任意维数的多重初等邻域，而初等结构态射的相对维数为1；下一段的归纳也确认这里应为“多重初等”。法稿句末“dans”后的原字受改写影响，不能可靠辨清；英文排成$X$并附“最后这个字母是否正确？字迹不清”的疑问。正文的“在$X$中”据上下文整理，不冒称辨清了原稿该字；英文疑问在此保存。}

对$X$的维数$n$作归纳，我们归约为证明：若$n>0$，则存在$x$的一个开邻域$U$以及一个初等态射$f\colon U\to V$，其中$V$是$k$上的光滑概形，必然几何不可约且维数为$n-1$。（$n=0$的情形显然平凡，因为此时$X$在$k$上同构于$\Spec(k)$；鉴于13.1并未排除空族态射的复合，它在$\Spec(k)$上是多重初等的。我们本应在13.1中以某种方式指出这一点。）\cmt{V21-INI005}{末句是原作者的编写提醒。零维结论在先取光滑有限型邻域后成立：该邻域是有限平展的，几何不可约使其几何纤维只有一点，因而就是$\Spec(k)$。归纳的实际拼接还要在$V$上取$f(x)$的多重初等邻域，再沿$f$拉回；初等态射对基变换稳定。一般闭点的有限集合版归纳由13.4(c)的补证处理。本注及报告中的展开均为译者核验，不是原稿省略文字的复原。}

% V21-INI-P03: French 74 bottom--75 top, dimension one.
首先，要求$k$完美的必要性，在$n=1$的情形就已显现：此时取$X$的函数域$K$的典范射影正规模型作为$X'$（参见第II章第7段）。\footnote{\textbf{英译编者注〔英文脚注64〕：}EGA II，是的，7.4〔英译者注〕。}$k$完美保证$X'$在$k$上光滑，因为$X'$无论如何都是正则的；它也保证，赋予诱导约化结构的$Z=X'\setminus X$在$k$上平展。现讨论一般情形，可假设$n\geq2$。\cmt{V21-INI006}{先取光滑仿射曲线邻域，可将它嵌入其正规射影模型；其补集非空，否则正维曲线将同时仿射且固有。补集是有限个闭点，完美性使其约化结构有限平展，因而满足定义13.1的满射及次数条件。原稿指出证明中使用了完美性，这一观察本身还不是不可删去假设的反例。报告“初段补证二”另给出$k=\mathbb F_3(t)$上正则而非光滑的射影三次曲线，证明在一般不完美域上结论确实可失败；该反例是译者补证，不是原文内容。}

% V21-INI-P04: French 75, compactification and normalization.
显然可以假设$X$仿射，因而拟射影。再把$X$替换为一个射影闭包，就可以假设$X$射影，但始终须证明：$x$的每个邻域都包含一个允许初等态射$U\to V$的开邻域$U$。同样，把$X$替换为其正规化（它在$X$上有限，因而射影；引用待补），\footnote{\textbf{英译编者注〔英文脚注65〕：}我想是在EGA II〔英译者注〕。}并不会改变$x$的一个邻域，所以可以假设$X$正规；由于$k$完美，它因而在$k$上几何正规。\cmt{V21-INI007}{法稿确实留下“r\'ef.”，英文脚注65仍带疑问，二者作为原有编写状态保留。数学上，在先取有限型光滑仿射邻域之后，射影闭包为$k$上的代数概形，其正规化有限且射影；正规化在原光滑邻域上为同构。完美域上的正规概形几何正规。报告“初段补证一”列出所用条件与可核验出处；这一补证没有把原稿未填的历史引文冒充为已经定位。}

% V21-INI-P05: French 75, divisor and singular locus.
这一假设的好处是：$X$中在$k$上不光滑的点所成的集合$Z$，余维至少为2。选择一个射影浸入$i\colon X\to\mathbf P^r$。取各个$\mathcal O_{\mathbf P^r}(d)$（$d>0$）中在$x$处不消失的截面的非零点开集，就得到$x$在$\mathbf P^r$中的一个基本邻域系。因此，$x$的每个邻域都包含一个形如$X\setminus Y$的邻域，其中$Y$是$X$中包含$Z$的闭子集，纯维数为$n-1$，且$x\notin Y$。\cmt{V21-INI008}{法稿将扭曲次数也记为$n>0$；为与$\dim X=n$区别，正文沿用旧译的$d$。“截面”所给的邻域明确为其非零点开集。选择$Y$的补证是：在给定开集的补集与$Z$的并上消失、而在$x$处不消失的某个正次数齐次式，给出$X$上的非零有效Cartier除子；取其约化支集即得所需纯余维1的$Y$。这里使用已经完成的有限型、整、射影归约，详见报告。}

赋予$Y$诱导约化结构；由于$k$完美，$Y$的奇异点集的维数至多为$n-2$。扩大前面的集合$Z$，便得到一个维数至多为$n-2$的闭子集$Z\subset Y$，它包含$X$和$Y$的几何奇异点集。

% V21-INI-P06: French 75 bottom--76 top, idea and lemma.
证明的想法是：固定一个余维为$n$的线性子簇$C$，用$X$与$P$中包含$C$的余维$n-1$线性子簇$L$的交，来把$X$纤维化。为此，需要如下引理。

\egas{13.3}{引理}沿用前述记号，设$X\supset Y\supset Z$是$\mathbf P^r_k$的闭子概形，维数分别为$n$、$n-1$及至多$n-2$，且$X\setminus Z$、$Y\setminus Z$光滑。$Z$的维数至多为$n-2$；必要时，若$k$的特征为$p>0$，按第9节的办法把射影浸入$i\colon X\to\mathbf P^r$换成任意一个“倍数”$\phi_d\circ i$（$d\geq2$），则存在$\mathbf P^r$中一个余维为$n-1$的线性子簇$L_0$，具有以下良好性质：\cmt{V21-INI009}{法稿第76页清楚写$X\supset Y\supset Z$、复合$\phi_n i$以及“quelconque”（任意）；英文的$X\supset Y\subset Z$、$\phi_n^i$与“some?”均属于英文转写或编辑问题。本版据法稿恢复，不再把它们归咎作者。正文把倍数指标$n$改成$d$，并明确写复合号。原稿重复的$Z$维数条件仍保留；原划去的逐点维数措辞没有重新补作有效正文。英文脚注66仅标“Fr”（法文），并非额外数学条件；该编辑标记在此保存。}
\begin{enumerate}[label=\alph*)]
\item $L_0\cap Z=\varnothing$；
\item $L_0\cap X$光滑，维数为1；
\item $L_0\cap Y$光滑，维数为0。
\end{enumerate}
（注意：这里$k$是无限域，并不要求它完美。）

% V21-INI-P07: French 76, start of construction.
暂且承认这个引理，并使用其后文证明中$L_0$经过$x$的形式，说明如何推出存在$x$的一个开邻域$U\subset X\setminus Y$，使得存在初等态射$U\to V$。\cmt{V21-INI010}{法稿的引理陈述没有写$x\in L_0$，但第78页证明的第一句明确要求截面经过$x$；末尾要让所得开集包含$x$也必须使用这个条件。正文在调用处补明，不把它冒称为引理原陈述已有条件。原证明在后面具体取$x\in P(k)$，因此下面按这一有理点构造理解；没有据此缩小13.2的原陈述。一般闭点不能直接当成$k$-有理点，所需有限几何点构造和下降另见本版13.4(c)的译者补证。}

存在$L_0$的一个线性子簇$C$，它在$\mathbf P^r$中余维为$n$，即在$L_0$中余维为1，不与有限集$(L_0\cap Y)\cup\{x\}$相交，并且横截光滑曲线$L_0\cap X$。令$T=X\cap C$；于是$T$是$X$的一个非空闭子概形，在$k$上平展，不含$x$且与$Y$不相交。\cmt{V21-INI011}{法稿明确写$C\subset L_0$，英文两处漏掉下标0，旧D17-M056对应的问题只应归于英文。法稿保留下来的选取句只明列避开有限集，随后却断言$T$平展；避开有限集本身不能推出横截性。原页被划去的旧句曾援引第7节并要求交光滑、零维，这属于作者删改层，不能重新当成有效原句。本版正文补明横截条件，译者在报告“初段补证三”证明满足横截与避开条件的超平面构成非空开集；$k$无限保证可在$k$上选取。正维射影曲线与超平面必相交，因此$T$非空。这项补证明确区别于对删划文字的辨读。}

% V21-INI-P08: French 76 bottom--77 top, parameter space and families.
另一方面，考虑$\operatorname{Grass}^{n-1}(P)$中对应于包含$C$的线性子簇的子概形$Q$。$Q$是一个维数为$n-1$的射影空间；特别地，它在$k$上光滑，维数为$n-1$。于是$L_0$对应于一个点$\xi_0\in Q(k)$。沿用此前引入的一般记号，考虑$X^{(n-1)}$沿浸入$Q\hookrightarrow\operatorname{Grass}^{n-1}(P)$的逆像$X_Q^{(n-1)}$，以及相应的逆像$Y_Q^{(n-1)}$和$T_Q^{(n-1)}$；后二者也是$X_Q^{(n-1)}$中互不相交的闭子概形。以$p,q,r$分别记这三个概形到$Q$的结构投影。\cmt{V21-INI012}{法稿第76页的参数点清楚为$\xi_0$，第77页的第三个逆像清楚为$T_Q^{n-1}$；英文前者漏下标0，后者在第50页末漏对象字母，旧D17-M057的纠错据此限定为英文问题。法稿定义$Q$时“包含”后的单个字母覆打不清，英文为$C$；按固定中心的定义正文写$C$，不声称该覆打字已完全辨明。上标加括号沿用现稿对余维参数的排法；法稿给关联族加的底线与当前统一的关联族记号对应，未另造新对象。}

由假设，$p$在$\xi_0$上方的各点光滑，$q$在$\xi_0$上方的各点平展；$r$亦如此，因为立即可见$T_Q^{(n-1)}$在$Q$上同构于$T\times_k Q$。\cmt{V21-INI013}{法稿把“由构造”删改为“按假设”。这里不应仅凭一个纤维光滑就断言态射光滑：$L_0$避开$Z$使$X,Y$在交点处光滑，且切出的维数分别恰为1和0，因此定义$L_0$的$n-1$个线性方程在相应切空间上的微分独立。相对Jacobian判据遂分别给出$p$光滑及$q$平展。$T\subset C$使每个参数截面都包含$T$，故第三个关联族确为乘积。完整核验见报告“初段补证四”。}

最后，$p$是固有态射；在这里$X$射影且几何不可约的条件下，$p$的纤维几何连通（Bertini定理）。\cmt{V21-INI014}{法稿本句只重提“$X$几何连通”，但前面已假设$X$几何不可约且射影。正文明确保留实际需要的完整条件。对任意几何参数，令其线性截面为$L$，则$\dim X+\dim L=n+(r-n+1)=r+1>r$；射影不可约簇的线性截面连通性定理因而适用。报告给出已核读的独立文献及其定理1.1，不再将旧D17-M058的“尚未核实”作为当前结论；同时不把单独的几何连通性误说成充分条件。}

% V21-INI-P09: French 77, restriction and elementary boundary.
因此，存在$\xi_0$在$Q$中的一个开邻域$V$，使得
\[
X'=X_Q^{(n-1)}\big|_V
\]
在$V$上固有且光滑，纤维几何连通。由于$\xi_0$处的纤维就是$X\cap L_0$，维数为1，缩小$V$后，还可假设$X'$在$V$上的纤维全都为一维。\cmt{V21-INI015}{法稿第77页正确写开集$V$及点$\xi_0$，但把环境写成$X$，应为参数空间$Q$；英文又误转成“$Q$在$X$中的开邻域$U$”。本版只把法稿真实错误归给作者，其余归于英文。缩小基底用到$p$固有：非光滑点集及不合所需相对维数的闭集，其像不含$\xi_0$，可从基底删去。实际上$p$射影，因为关联族闭嵌入$X\times_k Q$，而$X$射影；定义13.1所需的射影性并未缺失。}

最后，取$V$充分小，可假设$Y_Q^{(n-1)}|_V$及$T_Q^{(n-1)}|_V$在$V$上平展，于是二者的不交和预概形$Z'$（可识别为$X'$的一个闭子预概形）在$V$上平展。因此，令$U=X'\setminus Z'$，则$U\to V$是一个初等态射。\cmt{V21-INI016}{这里还须核验定义中的“满射、次数恒定”。$Z'$闭于射影族$X'$中，故在$V$上固有；结合平展性可知它有限平展。$V$是射影空间的非空开集，因而连通，有限平展次数遂恒定。其不交分量$T_Q^{(n-1)}|_V\simeq T\times_k V$非空且满射到$V$，故$Z'\to V$满射。报告将这些原稿略去的条件逐项补齐，没有改变定义。}

% V21-INI-P10: French 77 end; exact stop before French 78 lemma proof.
另一方面，$U$也是
\[
X''=X_Q^{(n-1)}\setminus
\bigl(Y_Q^{(n-1)}\cup T_Q^{(n-1)}\bigr)
\]
的一个开子集；$X''$是$X\setminus(Y\cup T)$在$X_Q^{(n-1)}$中的逆像。另一方面，态射$X''\to X\setminus(Y\cup T)$显然是同构，因为
\[
X_Q^{(n-1)}\setminus T_Q^{(n-1)}\longrightarrow X\setminus T
\]
是同构。\cmt{V21-INI017}{法、英两稿在$X''$的首个公式中都误写$Z_Q^{n-1}$；同一句立即说明这是$X-Y-T$的逆像，前文删去的边界也正是$Y_Q$与$T_Q$，故正文订为$T_Q^{(n-1)}$。这个错误确在法稿中。所用同构来自以$C$为中心的线性投影：对于$z\notin C$，存在唯一的余维$n-1$线性子簇$\langle C,z\rangle$；在概形上，中心外的关联族就是该投影的图，详见报告。}

因此，$U$可识别为$X\setminus(Y\cup T)$的一个开子集；它包含$(L_0\cap X)\setminus(Y\cup T)$，因而包含$x$。这就是所需的、包含在$X\setminus Y$中的$x$的邻域。\cmt{V21-INI018}{法、英两稿都说$U$包含整个$L_0\cap X$；这与已删去非空边界$T$相矛盾，故正文改为删去$Y\cup T$后的部分。包含$x$还用到调用时已补明的$x\in L_0$、原有的$x\notin Y$以及选取$C$时的$x\notin C$。本模块至此止；下一原段从证明引理13.3开始，不重复输入。}
